\documentclass[12pt, a4paper]{article}

\usepackage[a4paper, margin=2.5cm]{geometry}
\usepackage{setspace}
\usepackage{parskip}
\usepackage{booktabs}
\usepackage{array}
\usepackage{enumitem}
\usepackage{titlesec}
\usepackage{hyperref}
\usepackage{times}

\hypersetup{
    colorlinks=true,
    linkcolor=black,
    urlcolor=black,
    citecolor=black
}

\titleformat{\section}{\normalfont\large\bfseries}{\thesection}{1em}{}
\titlespacing*{\section}{0pt}{1.5ex plus 0.5ex minus 0.2ex}{0.8ex plus 0.2ex}

\setlength{\parindent}{0pt}
\setlength{\parskip}{8pt}

\begin{document}

% ── Title block ──────────────────────────────────────────────────────────────
\begin{center}
    {\Large \textbf{Automated Classification of Failed LLM-Generated\\[4pt] REST API Test Cases}}\\[10pt]
    {\normalsize \textit{Scope Description }}\\[6pt]
    {\normalsize Selen Öznur}\\[4pt]
    {\normalsize University of Antwerp}\\[4pt]
    {\normalsize July 2026}
\end{center}

\vspace{1em}

% ── 1. Context ────────────────────────────────────────────────────────────────
\section{Context}
REST APIs are the backbone of modern cloud-native systems. Testing them
comprehensively is challenging due to the vast input space, diverse endpoint
behaviors, and evolving specifications. Recent work has explored the use of
Large Language Model (LLM)-driven agents to automatically amplify existing
test suites.

A comparative study by Besjes et al.~\cite{besjes2025agentic}
evaluated both single-agent and multi-agent LLM workflows for REST API
test amplification across five cloud applications: Spotify, Google Drive,
PetStore, VAmPI, and Restful-Booker. While the results were promising
in terms of structural API coverage, a significant portion of the
generated test cases failed. These failures were manually analyzed and
categorized by the research team into five categories:

\begin{itemize}[leftmargin=1.5em, itemsep=3pt]
    \item \textbf{Semantically Incorrect}: the test logic is fundamentally inconsistent with the endpoint's intended behavior
    \item \textbf{Wrong Assertion}: the request is valid, but the expected outcome is incorrect
    \item \textbf{Missing Information}: essential information is unavailable in the test environment or API documentation
    \item \textbf{Wrong Input}: the test provides invalid, malformed, or incomplete input values
    \item \textbf{Bug}: the test exposes behavior that contradicts the intended API behavior as defined by the OpenAPI specification
\end{itemize}

This manual classification process produced a labeled dataset of failed test cases across multiple systems and agent configurations, forming the basis of this research.

% ── 2. Problem Statement ──────────────────────────────────────────────────────
\section{Problem Statement}

Manual classification of failed LLM-generated test cases is labor-intensive and does not scale. As agentic systems generate hundreds or thousands of test cases per run, the ability to automatically understand \textit{why} a test fails becomes critical for:

\begin{itemize}[leftmargin=1.5em, itemsep=3pt]
    \item Prioritizing which failures deserve developer attention
    \item Identifying systematic weaknesses in the LLM agent pipeline
    \item Enabling automated repair strategies targeted at specific failure types
\end{itemize}

Prior studies have primarily focused on LLM-based test generation,
test amplification, test repair, and the analysis of errors arising when
REST APIs are used as LLM tools
~\cite{wang2024softwaretesting,nooyens2025amplification,
bandlamudi2025framework}. The automated classification of failed
LLM-generated REST API test cases according to an existing manually
defined taxonomy remains comparatively underexplored. This study
investigates an approach to this classification problem.
% ── 3. Research Questions ─────────────────────────────────────────────────────
\section{Research Questions}

\textbf{RQ1:} To what extent can an LLM-based classifier automatically reproduce the manual failure categorization of the existing labeled dataset?

\textbf{RQ2:} How does the inclusion or exclusion of different input
features, such as test code, API specification, error message, and HTTP
method, affect failure classification performance?


\textbf{RQ3:} How does classification performance vary across different REST APIs and failure categories?

\textbf{RQ4:} How does failure classification performance vary across
selected LLMs under consistent prompting and evaluation conditions?
% ── 4. Approach ───────────────────────────────────────────────────────────────
\section{Approach}

The research will use the existing labeled dataset from the Test Amplification reproduction package as reference labels~\cite{testamplification}.
 The proposed approach consists of the following steps:
\begin{enumerate}[leftmargin=1.5em, itemsep=4pt]

    \item \textbf{Dataset analysis}: Explore the distribution of failure
    categories across systems and agent configurations. Identify class
    imbalance and inter-system variation.

 \item \textbf{Input preparation}: Extract and structure the available
information from the reproduction package for each failed test case,
including the test code, relevant API endpoint specification, execution
error message, and HTTP method. System and agent-configuration metadata
will be retained for API-level and configuration-level analysis, but will
not be used as direct classification inputs.

    \item \textbf{LLM-based classification}: Implement an automated
    classification pipeline that accesses two selected LLMs through their
    programming APIs. The selected model names, versions, prompts, and
    inference parameters will be documented to support reproducibility.
    Zero-shot, few-shot, and definition-guided prompting strategies will
    be evaluated to classify each failed test case into one of the five
    failure categories.~\cite{hannan2026fewshot}

    \item \textbf{Feature contribution analysis}: Conduct controlled
    feature ablation experiments to evaluate the contribution of each
    input source. A full-input configuration containing the test code,
    API specification, execution error message, and HTTP method will be
    compared with configurations in which one information source is
    removed at a time. The model, prompt template, dataset split, and
    inference parameters will be kept constant across configurations.
    Feature contribution will be estimated from the resulting change in
    macro-F1 and per-category F1-scores.
    \item \textbf{Cross-model comparison}: Compare GPT-4o mini and Claude
Haiku 4.5 using the same dataset, prompt templates, input configurations,
few-shot examples, and evaluation procedure. Fixed model versions will
be used and documented to support reproducibility. The comparison will
examine overall performance, per-category performance, and recurring
misclassification patterns across the two models.

  \item \textbf{Evaluation}: Use macro-F1 as the primary classification
metric because the failure categories may be imbalanced. Report accuracy,
balanced accuracy, Cohen's kappa, per-category precision, recall, and
F1-score as complementary measures. Use confusion matrices to identify
categories that are frequently confused. Compare feature configurations
and selected LLMs using changes in macro-F1 and per-category F1-scores,
and report results separately for each REST API. Finally, conduct a
qualitative analysis of misclassified cases to identify recurring error
patterns.

\end{enumerate}

% ── 5. Expected Contributions ─────────────────────────────────────────────────
\section{Expected Contributions}

\begin{enumerate}[leftmargin=1.5em, itemsep=4pt]

    \item A prototype and empirical evaluation of an automated approach
    for classifying failure modes in LLM-generated REST API test cases.

    \item An empirical assessment of the contribution of different input
    sources to failure classification performance.

    \item A comparison of selected LLMs under consistent prompting and
    evaluation conditions.

    \item Insights into recurring failure patterns in the studied agentic
    LLM test amplification workflows, which may guide future improvements
    to agent design.

\end{enumerate}

% ── 6. Relation to Existing Work ──────────────────────────────────────────────

\section{Relation to Existing Work}

This research extends the REST API test amplification study by
Besjes et al.~\cite{besjes2025agentic}, which produced and manually
labeled the dataset used in this work. The present study investigates
how this classification process can be automated.

Related work by Bandlamudi et al.~\cite{bandlamudi2025framework}
introduced an error taxonomy for REST APIs used as LLM tools. While
their focus is API readiness for agent-based tool use, this study
focuses on classifying failed LLM-generated REST API test cases.
The work is also situated within the broader area of LLM-based
software testing~\cite{wang2024softwaretesting}.

% ── 7. Preliminary Planning ───────────────────────────────────────────────────
\section{Preliminary Planning}

\vspace{0.5em}
\begin{tabular}{@{} c p{7cm} p{5cm} @{}}
    \toprule
    \textbf{Phase} & \textbf{Activity} & \textbf{Period} \\
    \midrule
    1 & Literature review + dataset familiarization & July 13--17 \\[4pt]
    2 & Dataset analysis + feature engineering      & July 17--31 \\[4pt]
    3 & LLM classifier design + experiments         & August 1--25 \\[4pt]
    4 & Evaluation + qualitative analysis           & August 25 -- September 5 \\[4pt]
    5 & research writing                              & Throughout / finalize by September 11 \\
    \bottomrule
\end{tabular}
\begin{thebibliography}{9}
\bibitem{besjes2025agentic}
J. Besjes, R. Nooyens, T. Bardakci, M. Beyazit, and S. Demeyer,
\href{https://arxiv.org/abs/2510.27417}
{\textit{Agentic LLMs for REST API Test Amplification:
A Comparative Study Across Cloud Applications}}.
arXiv preprint arXiv:2510.27417, 2025.

\bibitem{testamplification}
\href{https://doi.org/10.5281/zenodo.19251015}
{\textit{Reproduction Package -- A Comparative Study of Agentic LLMs
for REST API Test Amplification}}.
Zenodo, Version 1, 2026.

\bibitem{wang2024softwaretesting}
J. Wang, Y. Huang, C. Chen, Z. Liu, S. Wang, and Q. Wang,
\href{https://doi.org/10.1109/TSE.2024.3368208}
{\textit{Software Testing with Large Language Models:
Survey, Landscape, and Vision}}.
IEEE Transactions on Software Engineering,
vol. 50, no. 4, pp. 911--936, 2024.

\bibitem{nooyens2025amplification}
R. Nooyens, T. Bardakci, M. Beyazit, and S. Demeyer,
\href{https://arxiv.org/abs/2504.08113}
{\textit{Test Amplification for REST APIs via Single and
Multi-Agent LLM Systems}}.
International Conference on Testing Software and Systems,
pp. 161--177, 2025.

\bibitem{hannan2026fewshot}
M. A. Hannan, R. Ni, C. Zhang, L. Jia, R. Mangal, and C. S. Pasareanu,
\href{https://doi.org/10.14722/last-x.2026.23025}
{\textit{On the Difficulty of Selecting Few-Shot Examples for
Effective LLM-based Vulnerability Detection}}.
Workshop on LLM Assisted Security and Trust Exploration
(LAST-X), 2026.

\bibitem{bandlamudi2025framework}
J. Bandlamudi, R. Chaudhuri, N. Gantayat, S. Ghosh,
K. Mukherjee, P. Agarwal, R. Sindhgatta, and S. Mehta,
\href{https://arxiv.org/abs/2504.15546}
{\textit{A Framework for Testing and Adapting REST APIs as LLM Tools}}.
arXiv preprint arXiv:2504.15546, 2025.



\end{thebibliography}
\end{document}